\begin{textsample}{POS Dim 4 – human – Score 4.00 – es/es\_ef\_anos\_iniciais\_ciencias\_da\_natureza\_3\_p3.txt}  \label{ex:f4_pos_020}
Logo, nessa etapa do Ensino Fundamental, o ensino de Ciências se centrará não só no desenvolvimento das habilidades básicas ( observação, experimentação, descrição, identificação, descriminação, categorização, comparação, classificação, etc. ), como também no desenvolvimento das habilidades mediadoras na linguagem científica ( explicação, dedução, argumentação, diferenciação, analogia etc. ) ( CARVALHO, 2001 ).
Conforme descrito por Lorenzetti e Delizoicov ( 2001 ), o processo de alfabetização está etimologicamente \textbf{associado} ao processo de aquisição do alfabeto e das habilidades de \textbf{leitura} e \textbf{escrita}, enquanto o \textbf{letramento} está \textbf{associado} ao processo de aquisição das habilidades necessárias o uso que se faz da \textbf{leitura} e da \textbf{escrita} em seu contexto social.
E “ esta conceituação de \textbf{letramento}, transcendendo a de alfabetização, será de fundamental importância para o entendimento da alfabetização científica para as Séries Iniciais ” ( LORENZETTI e DELIZOICOV, 2001 ).
Ou seja, o que se espera para o Ensino de Ciências nos Anos Iniciais do Ensino Fundamental é exatamente esse caráter mais próximo do conceito de alfabetização, sem se abandonar as discussões já feitas sobre o uso de uma ou outra terminologia.

% matched lemmas: associar, escrita, leitura, letramento
\end{textsample}
